% Chapter 5.  Conclusions are stated only as strongly as the filled slots of
% Chapter 4 allow; every claim points back to the slot that holds its evidence.
\chapter[CONCLUSION AND RECOMMENDATIONS]{Conclusion and Recommendations}
\label{chap:conclusion}

This project set out to replace infrastructure-based navigation with a
vision-based pipeline that stays accurate when the warehouse changes, and to
decide between two ways of building it. At the reporting cutoff the pipeline
exists in both forms, the four experiments that judge it are defined for both
domains, and the four simulation slots hold results. The four real-robot slots
do not. The conclusions below are therefore conclusions about the simulated
warehouse, and each is bounded by the validity limits recorded beside it in
\cref{chap:results}.

\section{Conclusions by Experiment}
\label{sec:conclusions}

\subsection{Experiment 1: is vision earning its place?}

On five static recordings, a five-state wheel-plus-gyroscope EKF whose noise
model is calibrated per recording localises to \SIrange{0.19}{0.49}{\metre}
aligned ATE, ahead of stereo visual-inertial SLAM running on a hand-set
configuration on four of five recordings (\cref{tab:proprio-accuracy}). That
ordering is not a property of the architectures: giving the filter the visual
systems' own inertial noise model degrades it on every recording and puts
VINS-Fusion ahead on all five (\cref{sec:proprio-tuning}). The defensible
conclusion is narrower than either reading: over this range the noise model is
worth more than the estimator, and the visual systems have never been tuned to
this platform. Two results survive the control. The proprioceptive sensors fail
in complementary ways, the IMU holding heading and losing position, the wheels
the reverse, and the EKF is better than either in both quantities because it
observes the gyroscope bias rather than the heading
(\cref{sec:h1-falsified}). And every dead-reckoning covariance is overconfident
by about two orders of magnitude (\cref{tab:proprio-covariance}). Against the
project's target of \SI{0.05}{\metre} ATE, no estimator qualifies; against the
\SI{1.5}{\metre} RPE target at \SI{1}{\second}, all of them do in static
scenes. The comparison is confined to static, encoder-equipped recordings and
to a floor on which the wheels effectively cannot slip.

\subsection{Experiment 2: how good is the detector?}

The three-class detector is reliable at naming and unreliable at finding:
pooled precision 0.684, moving-object recall 0.591, conditional class F1 0.995
over 9\,860 frames (\cref{tab:yolo-performance,tab:yolo-classification}). It
misses about \SI{41}{\percent} of the moving objects that are scored, and
\SI{62}{\percent} if every occluded or tiny object is counted; that miss rate is
the ceiling on any filtering benefit downstream. The ground truth is projected
simulator geometry rather than annotation, which makes the boxes tighter and the
scored subset easier than a human-labelled set, and no average precision is
reported because a single confidence threshold was retained.

\subsection{Experiment 3: does filtering help?}

The question the project is named for has no valid answer yet. Sixty-six runs
exist, fully instrumented, but the VINS-Fusion filter never activated, eleven
of twelve ORB-SLAM3 treatment pairs received different input, and the scripted
dynamic worlds were found to be an invalid testbed, so every dynamic-scene
trajectory is withheld as a finding (\cref{sec:trajectory-results}). This is
untested, not tested and failed. The static runs do establish two things: the
two front ends reach medians of \SI{0.370}{\metre} and \SI{0.454}{\metre}
aligned ATE, and ORB-SLAM3's eleven-fold rotational penalty is produced by one
or two discrete events near inertial initialisation, not by its steady-state
tracking (\cref{sec:rpe}). Neither front end keeps up with the frame rate at the
replay rate used, and they hide that deficit differently, one by dropping
frames and one by queueing them (\cref{sec:block-performance}).

\subsection{Experiment 4: is the map usable?}

Where tracking is sound, the mapping pipeline is: the tiny-warehouse map is
\SI{96.6}{\percent} precise in what it marks occupied, recovers
\SI{89.5}{\percent} of the surfaces it observed, and contains no blind cells and
\SI{0.7}{\percent} proven false obstacles along an independently driven route
(\cref{tab:map-geometry,tab:map-drivability}). The whole-building recall of
\SI{67.2}{\percent} is route coverage, not mapping failure. The drivability half
of that result is dated, because its independent route is lost. On the large
warehouse the attempt failed through tracking loss in repetitive structure,
before the back end was exercised, and its artifacts were not retained.

\section{Limitations and Threats to Validity}
\label{sec:limitations}

The limitations of \cref{tab:limitations} carry over without softening. The
ones that bound the conclusions above most tightly are these. Every result is
from one laptop at one replay rate, so nothing here is evidence of real-time
capability on the NRU-51V. The proprioceptive baseline exists only for static,
encoder-equipped, near-slip-free recordings, with $n=5$ across four worlds and
no condition repeated more than twice. ORB-SLAM3's run-to-run spread is
unmeasured and its loop-closure configuration is absent from the run manifest,
so its loop-closure results cannot be attributed. The detector is scored against
projected geometry with exemptions the training convention lacks. The dynamic
worlds move props kinematically and so do not pose the problem the project
poses. And no real-robot slot is filled: the rig's calibration has a
placeholder translation and a reused intrinsic set, and no recording carries
motion-capture truth, wheel odometry or varied illumination.

\section{Recommendations and Next Steps}
\label{sec:recommendations}

\textbf{Fill the four real-robot slots, in the order the simulation results
suggest.} First calibrate the real camera--IMU pair and record the
motion-capture reference with the wheel and steering topics bridged, which
unlocks Experiments~1 and~3 at once and yields the first measured noise model of
the real IMU. Then annotate a frame subset under the three illumination levels
for Experiment~2, and survey the site into a reference grid for Experiment~4.
The integrated live test of \cref{sec:vio-filter-method}, both pipelines at
\SI{1.5}{\metre\per\second} under scheduled lighting change with the dynamic
variable triggered at each zone transfer, is the run that turns the eight-slot
table into a comparison.

\textbf{Give Experiment~3 a valid test before repeating it.} Repair the
detection-to-VINS integration and prove it by the removed-track counter;
input-match the ORB-SLAM3 treatment pairs by reducing filtered-mode latency;
rebuild the dynamic worlds with physically simulated props; and record the
\repo{allsensor} topics in a dynamic world so that the proprioceptive baseline
extends to the condition that matters.

\textbf{Tune the visual systems to the platform.} The cheapest remaining
improvement is a measured inertial noise model for the visual estimators; the
one attempt so far was bimodal, better on four of five recordings and divergent
on three of ten runs, and needs a proper stationary recording behind it.

\textbf{Record every behaviour-changing switch in the run manifest}, repeat
every nondeterministic configuration at least three times, and retain the
artifacts of failed attempts; the large-map failure and the lost drivability
route are both results the project cannot presently audit.

\textbf{Future direction.} Integrated relocalization against a saved map,
global pose-graph optimisation with loop-closure constraints, and alignment and
merging of maps across sessions remain future work, to be taken up only once the
front end tracks reliably in repeated structure; path planning stays outside the
project.
